\documentclass[aspectratio=169,11pt]{beamer}
\def\courseweek{9}
% Shared style for the Introduction to Information Geometry lecture slides.
% Compile the lecture .tex files from this directory.
\providecommand{\courseweek}{1}
\usepackage[T1]{fontenc}
\usepackage[utf8]{inputenc}
\usepackage[sc]{mathpazo}
\renewcommand{\sfdefault}{lmss}
\renewcommand{\ttdefault}{lmtt}
\usepackage{amsmath,amssymb,mathtools}
\usepackage{tikz}
\usetikzlibrary{arrows.meta,calc,positioning}
\usepackage{pgfplots}
\pgfplotsset{compat=1.18}
\usepackage{booktabs}

\definecolor{courseink}{HTML}{173749}
\definecolor{coursemuted}{HTML}{526674}
\definecolor{courseteal}{HTML}{0E7287}
\definecolor{coursecopper}{HTML}{B66B43}
\definecolor{courseline}{HTML}{D4E0E3}
\definecolor{coursepale}{HTML}{EFF5F5}

\setbeamercolor{normal text}{fg=courseink,bg=white}
\setbeamercolor{structure}{fg=courseteal}
\setbeamercolor{frametitle}{fg=courseink,bg=white}
\setbeamercolor{title}{fg=courseink}
\setbeamercolor{subtitle}{fg=coursemuted}
\setbeamercolor{block title}{fg=courseink,bg=white}
\setbeamercolor{block body}{fg=courseink,bg=white}
\setbeamercolor{alerted text}{fg=coursecopper}
\setbeamerfont{title}{family=\rmfamily,series=\bfseries,size=\fontsize{32}{37}}
\setbeamerfont{frametitle}{family=\rmfamily,series=\bfseries,size=\fontsize{26}{31}}
\setbeamerfont{normal text}{size=\fontsize{18}{22}}
\setbeamerfont{subtitle}{size=\fontsize{17}{21}}
\setbeamersize{text margin left=0.78cm,text margin right=0.78cm}
\setbeamertemplate{navigation symbols}{}
\setbeamertemplate{itemize items}[circle]
\setbeamertemplate{enumerate items}[default]
\setbeamertemplate{blocks}[default]
\setlength{\parskip}{0.22em}
\setbeamertemplate{footline}{%
  \leavevmode\hbox to \paperwidth{\hspace{.78cm}%
  {\color{coursemuted}\fontsize{8.5}{10}\selectfont Introduction to Information Geometry \hspace{.32cm} / \hspace{.32cm} Week \courseweek}%
  \hfill{\color{coursemuted}\fontsize{8.5}{10}\selectfont \insertframenumber\hspace{.17cm}/\hspace{.17cm}\inserttotalframenumber}\hspace{.78cm}}%
  \vskip.17cm
}
\setbeamertemplate{frametitle}{%
  \vspace{.13cm}
  \insertframetitle\par
  \vspace{.08cm}
}
\setbeamertemplate{title page}{%
  \vspace*{1.00cm}
  {\color{courseteal}\fontsize{10}{13}\selectfont\bfseries
  INTRODUCTION TO INFORMATION GEOMETRY\par}
  \vspace{.40cm}
  {\usebeamerfont{title}\usebeamercolor[fg]{title}\inserttitle\par}
  \vspace{.22cm}
  {\usebeamerfont{subtitle}\usebeamercolor[fg]{subtitle}\insertsubtitle\par}
  \vfill
  {\color{coursemuted}\fontsize{13}{16}\selectfont Christopher Lee\par}
  \vspace{.34cm}
}

\pgfplotsset{
  courseplot/.style={
    width=11.05cm,
    height=5.15cm,
    axis line style={coursemuted!75},
    tick style={coursemuted!65},
    ticklabel style={font=\scriptsize,color=coursemuted},
    label style={font=\normalsize,color=courseink},
    grid=major,
    major grid style={courseline!70,line width=.4pt},
    line width=1.45pt,
    legend style={font=\small,draw=none,fill=none,at={(.98,.03)},anchor=south east},
    clip=true
  }
}

\newcommand{\highlightmath}[1]{{\color{courseteal}#1}}
\newcommand{\thinrule}{\par\vspace{.25cm}{\color{courseline}\rule{\linewidth}{.8pt}}\par}

% A one-dimensional state space.  x=0 and x=1 are placed at 0 and 9 units.
\newcommand{\stateinterval}[1]{%
 \begin{tikzpicture}[x=1cm,y=1cm]
   \draw[courseink,line width=1.25pt] (0,0)--(9,0);
   \foreach \pos/\lab in {0/0,9/1} {
     \fill[courseink] (\pos,0) circle(2.7pt);
     \node[below=9pt,text=coursemuted] at (\pos,0) {\(\lab\)};
   }
   #1
 \end{tikzpicture}%
}

% Paired pages are separate frames in the PDF, with matching titles.
% Supply the whole frame body for each version; the instructor may edit them
% independently without changing overlay logic.
\newcommand{\coursepair}[3]{%
  \begin{frame}{#1}#2\end{frame}%
  \begin{frame}{#1}#3\end{frame}%
}

% Equilateral representation of the categorical simplex, in barycentric
% coordinates with vertices A=(0,0), B=(1,0), C=(1/2,sqrt(3)/2).
\newcommand{\coursesimplex}[1]{%
 \begin{tikzpicture}[scale=4.9]
  \fill[coursepale] (0,0)--(1,0)--(.5,.8660254)--cycle;
  \draw[courseink,line width=.65pt] (0,0)--(1,0)--(.5,.8660254)--cycle;
  \node[below left,font=\small] at (0,0) {A};
  \node[below right,font=\small] at (1,0) {B};
  \node[above,font=\small] at (.5,.8660254) {C};
  #1
 \end{tikzpicture}%
}

\title{Potentials, Stability, and KL Flows}
\subtitle{Week 9 \(\cdot\) Lecture 2}
\date{}
\begin{document}
\begin{frame}[plain]\titlepage\end{frame}

\begin{frame}{When a payoff matrix has a potential}
Suppose fitness is \(f(p)=Mp\), with \(M\) symmetric.
Define
\[
\highlightmath{\Phi(p)=\tfrac12p^TMp}.
\]
\vspace{.47cm}
Then \(\partial_i\Phi=(Mp)_i=f_i\), and the replicator
equation is the Fisher gradient of \(\Phi\).
\end{frame}

\begin{frame}{An example that favors balance}
Choose \(M=-I\). Each strategy's fitness is \(f_i=-p_i\), and
\[
\highlightmath{\Phi(p)=-\frac12\sum_i p_i^2}.
\]
\vspace{.43cm}
For a fixed total probability, the sum of squares is smallest
at the uniform distribution.
\vspace{.35cm}
Thus the uniform state maximizes this potential.
\end{frame}

\coursepair{Compute a balancing flow}{
For \(M=-I\), calculate the replicator velocity at
\[
P=(.6,.3,.1).
\]
\vspace{.32cm}
Use \(f_i=-p_i\) and \(p_i'=p_i(f_i-\bar f)\).
}{
For \(P=(.6,.3,.1)\), the mean fitness is
\(\bar f=-\sum p_i^2=-.46\).
\vspace{.35cm}
Therefore
\[
\highlightmath{p'=(-.084,\ .048,\ .036)}.
\]
\vspace{.24cm}
The largest probability falls; the smaller two rise.
}

\coursepair{Verify ascent at the same state}{
At \(P=(.6,.3,.1)\), we found
\(p'=(-.084,.048,.036)\).
\vspace{.39cm}
Evaluate the potential derivative using
\[
\dot\Phi=\sum_i f_i p_i'=-\sum_i p_i p_i'.
\]
}{
At \(P=(.6,.3,.1)\),
\[
\dot\Phi=-(.6)(-.084)-(.3)(.048)-(.1)(.036).
\]
\vspace{.28cm}
Thus
\[
\highlightmath{\dot\Phi=.0324>0}.
\]
}

\begin{frame}{Curvature of the potential and stability}
At an interior critical point \(q\), the derivative of a metric
gradient field is
\[
\highlightmath{DF(q)=G(q)^{-1}\operatorname{Hess}\Phi(q)}.
\]
\vspace{.51cm}
If the Hessian is negative definite on the tangent space,
then \(q\) is locally attracting.
\end{frame}

\begin{frame}{Why gradient equilibria have real eigenvalues}
The linearization \(G^{-1}H\) need not itself be symmetric.
\vspace{.45cm}
But it is similar to
\[
\highlightmath{G^{-1/2}HG^{-1/2}},
\]
\vspace{.40cm}
which is symmetric whenever \(H=\operatorname{Hess}\Phi\) is symmetric.
\vspace{.36cm}
Consequently its eigenvalues are real.
\end{frame}

\begin{frame}{A target distribution as a potential}
Let \(Q=(q_1,\ldots,q_n)\) have strictly positive coordinates.
\vspace{.32cm}
Define
\[
\highlightmath{\Phi(P)=-D(Q\Vert P)
=-\sum_i q_i\ln\frac{q_i}{p_i}}.
\]
\vspace{.38cm}
The potential is largest when \(P=Q\).
What does its Fisher-gradient flow look like?
\end{frame}

\begin{frame}{The target-seeking flow is linear}
Differentiate the potential with respect to \(p_i\):
\[
\frac{\partial\Phi}{\partial p_i}=\frac{q_i}{p_i},
\qquad\sum_i p_i\frac{\partial\Phi}{\partial p_i}=1.
\]
\vspace{.39cm}
The natural-gradient equation simplifies to
\[
\highlightmath{p_i'=q_i-p_i}.
\]
\vspace{.34cm}
Every coordinate relaxes toward its specified target.
\end{frame}

\coursepair{Solve the target-seeking system}{
Consider \(P'=Q-P\), where \(Q\) is fixed.
\vspace{.49cm}
Solve the initial-value problem \(P(0)=P_0\).
\vspace{.34cm}
What happens to the distance \(P(t)-Q\) as time increases?
}{
Consider \(P'=Q-P\), with \(P(0)=P_0\).
\vspace{.47cm}
Each coordinate satisfies a first-order linear equation:
\[
\highlightmath{P(t)=Q+e^{-t}(P_0-Q)}.
\]
\vspace{.29cm}
The displacement from the target decays by \(e^{-t}\).
}

\coursepair{One step toward a chosen target}{
Let
\[
Q=(.3,.4,.3),\qquad P_0=(.65,.2,.15).
\]
\vspace{.38cm}
Use \(P(1)=Q+e^{-1}(P_0-Q)\).
\vspace{.32cm}
Find \(P(1)\) to three decimal places.
}{
Let \(Q=(.3,.4,.3)\) and \(P_0=(.65,.2,.15)\).
\[
P(1)=(.3+.35/e,\ .4-.2/e,\ .3-.15/e).
\]
\vspace{.24cm}
Thus
\[
\highlightmath{P(1)\approx(.429,.326,.245)}.
\]
}

\begin{frame}{Reverse the KL divergence}
Now choose the distinct potential
\[
\Phi(P)=-D(P\Vert Q)=-\sum_i p_i\ln\frac{p_i}{q_i}.
\]
\vspace{.43cm}
Its Fisher gradient is
\[
\highlightmath{p_i'=p_i\left[-\ln\frac{p_i}{q_i}
+D(P\Vert Q)\right]}.
\]
\end{frame}

\begin{frame}{Reverse-KL flow has a geometric interpolation}
Let \(a(t)=e^{-t}\). The reverse-KL solution is
\[
\highlightmath{p_i(t)=\frac{q_i^{1-a(t)}p_i(0)^{a(t)}}
{\sum_j q_j^{1-a(t)}p_j(0)^{a(t)}}}.
\]
\vspace{.41cm}
This path interpolates multiplicatively between the initial
distribution and the target, then renormalizes.
\end{frame}

\coursepair{Reverse-KL flow toward uniformity}{
Suppose \(Q=(1/3,1/3,1/3)\) and
\(P_0=(.64,.25,.11)\).
\vspace{.42cm}
At \(t=\ln2\), find the unnormalized weights
\(p_i(0)^{e^{-t}}\).
\vspace{.29cm}
How would you turn them into a probability distribution?
}{
Suppose \(Q=(1/3,1/3,1/3)\) and
\(P_0=(.64,.25,.11)\).
\vspace{.33cm}
At \(t=\ln2\), \(e^{-t}=1/2\), so the weights are
\[
\highlightmath{(.8,.5,\sqrt{.11})}.
\]
\vspace{.26cm}
Dividing by \(.8+.5+\sqrt{.11}\) gives
\(P(t)\approx\highlightmath{(.490,.306,.203)}\).
}

\begin{frame}{Why every flow is not a gradient flow}
An increasing potential cannot return to its starting value
along a nonstationary closed orbit.
\vspace{.55cm}
The balanced three-strategy cyclic replicator has closed orbits.
\vspace{.52cm}
Such cyclic motion cannot be globally realized as ascent of a
single smooth Fisher potential.
\end{frame}

\begin{frame}{Dynamics constrain the possible geometry}
When fitness comes from a potential, its Fisher-gradient flow
provides monotone ascent and a direct stability test.
\vspace{.48cm}
When motion cycles, a single global ascent potential is unavailable.
\vspace{.48cm}
These two behaviors give us complementary ways to organize
the dynamics of probability distributions.
\end{frame}
\end{document}
